\begin{corollary}{2.2.1}
\label{II.2.2.1}
Let \(X\) be an \(S\)-scheme and \(\mathcal{M}\) a free \(\OO_S\)-module of finite
type.
The functor \(\Hom_S(\mathrm{I}_S(\mathcal{M}),X)\) is isomorphic (as a functor
over \(X\)) to a finite product (over \(X\)) of copies of
\(\Hom_S(\mathrm{I}_S,X)\).
\end{corollary}

\begin{enoncerm}{Nota}{2.2.2}
\label{II.2.2.2}
It follows from the proof of the proposition that
\(\Hom_S(\mathrm{I}_S,X)\) is isomorphic as an \(X\)-functor to
\(\mathbf{V}(\Omega^1_{X/S})\) (I~4.6.1), hence representable by the vector
fibration%
{\renewcommand{\thefootnote}{(12)}\nde{Cf.\ N.D.E.~(33) of Expos\'e~I.}}
\(\mathbf{V}(\Omega^1_{X/S})\).
\end{enoncerm}

\sgasection{3}{The tangent bundle, the condition (E)}
\label{II.sec.3}

In this section, unless the contrary is indicated, the letters \(\mathcal{M}\),
\(\mathcal{M}'\), \(\mathcal{N}\), etc., will always denote free
\(\OO_S\)-modules of finite type (that is, isomorphic to a finite direct sum of
copies of \(\OO_S\)).

We shall systematically use the identifications justified in Expos\'e~I; thus
we shall say ``functor over \(S\)'' to designate indiscriminately a functor
equipped with a morphism to \(S\) (\(=\mathbf{h}_S\)) or a functor on the
category of objects over \(S\). One will likewise say ``group functor over
\(S\)'', etc.

\begin{definition}{3.1}
\label{II.3.1}
Let \(S\) be a scheme and \(\mathcal{M}\) a free \(\OO_S\)-module of finite
type. Let \(X\) be a functor over \(S\). One calls the tangent bundle to \(X\)
over \(S\) relative to the \(\OO_S\)-module \(\mathcal{M}\), and one denotes by
\(\mathrm{T}_{X/S}(\mathcal{M})\), the \(S\)-functor
\[
\mathrm{T}_{X/S}(\mathcal{M})=\Hom_S(\mathrm{I}_S(\mathcal{M}),X).
\]
In particular, one calls the tangent bundle to \(X\) over \(S\), and one
denotes by \(\mathrm{T}_{X/S}\), the functor
\[
\mathrm{T}_{X/S}=\mathrm{T}_{X/S}(\OO_S)=\Hom_S(\mathrm{I}_S,X).
\]%
{\renewcommand{\thefootnote}{(13)}\nde{When \(S=\Spec(k)\) and \(X\) is a
\(k\)-scheme, one has
\(\mathrm{T}_{X/S}(S')=\Hom_S(S'\otimes_k k[\varepsilon],X)
=X(S'\otimes_k k[\varepsilon])\); one thus recovers one of the usual
definitions of the tangent bundle.}}
\end{definition}

Then \(\mathcal{M}\mapsto\mathrm{T}_{X/S}(\mathcal{M})\) is a covariant functor
from the category of free \(\OO_S\)-modules of finite type to the category of
\(S\)-functors. In particular the morphisms%
{\renewcommand{\thefootnote}{(14)}\nde{``\(0\to\mathcal{M}\) and
\(\mathcal{M}\to 0\)'' has been corrected to: ``\(\mathcal{M}\to 0\) and
\(0\to\mathcal{M}\)''.}}
\(\mathcal{M}\to 0\) and \(0\to\mathcal{M}\) define respectively an
\(S\)-morphism
\[
\pi_{\mathcal{M}}\colon
\mathrm{T}_{X/S}(\mathcal{M})\to\mathrm{T}_{X/S}(0)\simeq X
\]
and a section \(\tau_0\) of the latter, called the zero section (or null
section).

\begin{remark}{3.1.1}
\label{II.3.1.1}
{\renewcommand{\thefootnote}{(15)}\nde{Paragraphs 3.1.1 and 3.1.2 have been
added.}}
One will note that the projection
\(\pi_{\mathcal{M}}\colon\mathrm{T}_{X/S}(\mathcal{M})\to X\) is induced by
the zero section \(\varepsilon_{\mathcal{M}}\colon S\to\mathrm{I}_S(\mathcal{M})\),
while the null section
\(\tau_0\colon X\to\mathrm{T}_{X/S}(\mathcal{M})\) is induced by the structural
morphism \(\rho\colon\mathrm{I}_S(\mathcal{M})\to S\); \ie for every point
\(t\in\mathrm{T}_{X/S}(\mathcal{M})(S')\)
(\resp \(x\in X(S')\)), corresponding to an \(S\)-morphism
\(f\colon S'\times_S\mathrm{I}_S(\mathcal{M})\to X\)
(\resp \(g\colon S'\to X\)), one has
\(\pi(t)=f\circ(\mathrm{id}_{S'}\times\varepsilon_{\mathcal{M}})\)
(\resp \(\tau_0(x)=g\circ(\mathrm{id}_{S'}\times\rho)\)).
\end{remark}

It follows from 3.1 that \(\mathcal{M}\mapsto\mathrm{T}_{X/S}(\mathcal{M})\)
is a covariant functor from the category of free \(\OO_S\)-modules of finite
type to that of functors over \(X\). In particular \(\mathbf{O}(S)\) is a
monoid of operators of the \(X\)-functor \(\mathrm{T}_{X/S}(\mathcal{M})\)
which respects ``the functoriality in \(\mathcal{M}\)''.

\begin{scholium}{3.1.2}
\label{II.3.1.2}
\textup{\textsuperscript{(15)}}
What precedes means, in particular, the following things. For every
\(S\)-morphism \(\phi\colon X'\to X\), set
\[
\Sigma(X',\mathcal{M})=\Hom_X(X',\mathrm{T}_{X/S}(\mathcal{M})).
\]
One has an action of the multiplicative monoid
\(\End_{\OO_S}(\mathcal{M})\) on \(\Sigma(X',\mathcal{M})\), denoted
\((\lambda,x)\mapsto\lambda\ast x\), such that
\(\lambda\ast(\mu\ast x)=(\lambda\mu)\ast x\), \(1\ast x=x\), and
\(0\ast x=\tau_0\circ\phi\), where \(\tau_0\) is the null section
\(X\to\mathrm{T}_{X/S}(\mathcal{M})\). One likewise has an action of
\(\End_{\OO_S}(\mathcal{M}\oplus\mathcal{M})\) on
\(\Sigma(X',\mathcal{M}\oplus\mathcal{M})\).%
{\renewcommand{\thefootnote}{(16)}\nde{In fact, one is interested only in the
action of \(\mathbf{O}(S)\) (\resp \(\mathbf{O}(S)\times\mathbf{O}(S)\)) on
\(\Sigma(X',\mathcal{M})\)
(\resp \(\Sigma(X',\mathcal{M}\oplus\mathcal{M})\)), \cf below and the proof
of~3.6.}}

Moreover, let \(m\colon\mathcal{M}\oplus\mathcal{M}\to\mathcal{M}\)
(\resp \(\delta\colon\mathcal{M}\to\mathcal{M}\oplus\mathcal{M}\)) be the
addition (\resp the diagonal map) of \(\mathcal{M}\), and denote by
\(m_{X'}\colon\Sigma(X',\mathcal{M}\oplus\mathcal{M})\to\Sigma(X',\mathcal{M})\)
and
\(\delta_{X'}\colon\Sigma(X',\mathcal{M})\to\Sigma(X',\mathcal{M}\oplus\mathcal{M})\)
the morphisms induced by \(m\) and \(\delta\). For \(\lambda,\mu\in\mathbf{O}(S)\),
denote by \(\ell_{\lambda}\) (\resp \(\ell_{\lambda,\mu}\)) multiplication by
\(\lambda\) in \(\mathcal{M}\) (\resp by \((\lambda,\mu)\) in
\(\mathcal{M}\oplus\mathcal{M}\)).
Since \(m\circ\ell_{\lambda,\lambda}=\ell_{\lambda}\circ m\) and
\(m\circ\ell_{\lambda,\mu}\circ\delta=\ell_{\lambda+\mu}\), one has, for every
\(z\in\Sigma(X',\mathcal{M}\oplus\mathcal{M})\) and
\(x\in\Sigma(X',\mathcal{M})\):
\[
(\dagger)\qquad
\lambda\ast m(z)=m((\lambda,\lambda)\ast z),\qquad
m((\lambda,\mu)\ast\delta(x))=(\lambda+\mu)\ast x.
\]
\end{scholium}

\begin{definition}{3.2}
\label{II.3.2}
Let \(u\in X(S)=\Hom_S(S,X)=\Gamma(X/S)\). One calls the tangent space to
\(X\) over \(S\) at the point \(u\) relative to \(\mathcal{M}\), and one denotes
by \(\mathrm{L}^{u}_{X/S}(\mathcal{M})\), the \(S\)-functor obtained from the
\(X\)-functor \(\mathrm{T}_{X/S}(\mathcal{M})\) by inverse image under the
morphism \(u\colon S\to X\):
\[
\xymatrix{
  \mathrm{L}^{u}_{X/S}(\mathcal{M}) \ar[r] \ar[d]
    & \mathrm{T}_{X/S}(\mathcal{M}) \ar[d]^{\pi} \\
  S \ar[r]^{u}
    & X.
}
\]
In particular \(\mathrm{L}^{u}_{X/S}(\OO_S)\) is denoted by
\(\mathrm{L}^{u}_{X/S}\). It is the tangent space to \(X\) over \(S\) at the
point \(u\).
\end{definition}

\begin{remark}{3.2.1}
\label{II.3.2.1}
{\renewcommand{\thefootnote}{(17)}\nde{This remark has been added.}}
It follows from 3.1.1 that, for every \(t\colon S'\to S\),
\(\mathrm{L}^{u}_{X/S}(\mathcal{M})(S')\) is the set of \(S\)-morphisms
\(f\colon\mathrm{I}_{S'}(\mathcal{M})\to X\) such that
\(f\circ\varepsilon_{\mathcal{M}}=u\circ t\), where
\(\varepsilon_{\mathcal{M}}\colon S'\to\mathrm{I}_{S'}(\mathcal{M})\) is the
zero section.
\end{remark}

Let us immediately record the

\begin{proposition}{3.3}
\label{II.3.3}
If \(X\) is representable by an \(S\)-scheme denoted \(\widetilde{X}\), then
\(\mathrm{T}_{X/S}(\mathcal{M})\) and \(\mathrm{L}^{u}_{X/S}(\mathcal{M})\) are
representable. In particular \(\mathrm{T}_{X/S}\) and
\(\mathrm{L}^{u}_{X/S}\) are representable by vector fibrations over
\(\widetilde{X}\) and over \(S\) which are respectively
\(\mathbf{V}(\Omega^1_{\widetilde{X}/S})\) and
\(\mathbf{V}(u^*(\Omega^1_{\widetilde{X}/S}))\).
\end{proposition}

It obviously suffices to prove the proposition for
\(\mathrm{T}_{X/S}(\mathcal{M})\), the analogous results for
\(\mathrm{L}^{u}_{X/S}(\mathcal{M})\) being deduced from it by inverse image.
By Corollary~2.2.1, it even suffices to do so for \(\mathrm{T}_{X/S}\), and
in this case the proposition is nothing other than the remark indicated
in~2.2.2.

\begin{remark}{3.3.1}
\label{II.3.3.1}
From this proposition there follows a particularly simple description of the
vector fibration representing \(\mathrm{L}^{u}_{X/S}\): the image of the
section \(u\) of \(X\) over \(S\) is locally closed%
{\renewcommand{\thefootnote}{(18)}\nde{\Cf EGA~I, 5.3.11.}},
hence defined by a quasi-coherent ideal \(\mathfrak{m}\) of a scheme induced
on an open of \(X\). The quotient \(\mathfrak{m}/\mathfrak{m}^2\) may be
regarded as a quasi-coherent module over \(S\). It is this module which
defines the vector fibration sought.

For example, let \(X\) be an algebraic scheme over a field \(k\) and \(u\) a
point of \(X\) rational over \(k\). Let \(\mathfrak{m}\) be the maximal ideal
of the local ring \(\OO_{X,u}\) and let \(t\) be the \(k\)-vector space dual
of \(\mathfrak{m}/\mathfrak{m}^2\); it is the Zariski tangent space of
\(\OO_{X,u}\) at the point \(u\). Then, with the notations of I~4.6.5.1, one
has:
\[
\mathrm{L}^{u}_{X/k}=\mathbf{V}(\mathfrak{m}/\mathfrak{m}^2)=\mathbf{W}(t).
\]

This parenthesis being closed, let us return to the general situation. Let us
first remark that \(\mathrm{L}^{u}_{X/S}(\mathcal{M})\) is a covariant functor
from the category of free \(\OO_S\)-modules of finite type to that of functors
over \(S\). In particular \(\mathbf{O}(S)\) is a set of operators of the
\(S\)-functor \(\mathrm{L}^{u}_{X/S}(\mathcal{M})\) which respects the
functoriality in \(\mathcal{M}\).%
{\renewcommand{\thefootnote}{(19)}\nde{\Cf 3.1.2.}}
\end{remark}

\begin{proposition}{3.4}
\label{II.3.4}
% typo?: statement prints L_{X/S}(M) with no superscript u; kept as printed
The formation of \(\mathrm{T}_{X/S}(\mathcal{M})\) and
\(\mathrm{L}_{X/S}(\mathcal{M})\) commutes with extension of the base: for
every \(S\)-scheme \(S'\), one has isomorphisms functorial in \(\mathcal{M}\)
\begin{align*}
\mathrm{T}_{X_{S'}/S'}(\mathcal{M}\otimes\OO_{S'})
  &\isomto\mathrm{T}_{X/S}(\mathcal{M})_{S'},\\
\mathrm{L}^{u'}_{X_{S'}/S'}(\mathcal{M}\otimes\OO_{S'})
  &\isomto\mathrm{L}^{u}_{X/S}(\mathcal{M})_{S'},
  \qquad\text{where }u'=u_{S'}.
\end{align*}
\end{proposition}

This follows immediately from the fact that the \(\Hom\) commute with
extension of the base.

\begin{corollary}{3.4.1}
\label{II.3.4.1}
The \(X\)-functor \(\mathrm{T}_{X/S}(\mathcal{M})\) (\resp the \(S\)-functor
\(\mathrm{L}^{u}_{X/S}(\mathcal{M})\)) is endowed naturally with a structure of
object with operators \(\mathbf{O}_X\) (\resp \(\mathbf{O}_S\)), this
structure being functorial in \(\mathcal{M}\).
\end{corollary}

Let us first show this for \(\mathrm{L}^{u}_{X/S}(\mathcal{M})\). For each
\(S'\) over \(S\), \(\mathbf{O}(S')\) operates on
\(\mathcal{M}\otimes\OO_{S'}\), hence on
\(\mathrm{L}^{u'}_{X_{S'}/S'}(\mathcal{M}\otimes\OO_{S'})
=\mathrm{L}^{u}_{X/S}(\mathcal{M})_{S'}\); now one checks that this operation
is functorial in \(S'\). It therefore endows
\(\mathrm{L}^{u}_{X/S}(\mathcal{M})\), as announced, with a structure of
functor with operators \(\mathbf{O}_S\).

For \(\mathrm{T}_{X/S}(\mathcal{M})\) it is a little more complicated. For each
\(X'\) over \(X\), set
\(\mathrm{T}_{X/S}(\mathcal{M})_{X'}
=\mathrm{T}_{X/S}(\mathcal{M})\times_X X'\);
one must endow
\(\mathrm{T}_{X/S}(\mathcal{M})_{X'}(X')
=\Hom_X(X',\mathrm{T}_{X/S}(\mathcal{M}))\)
with a structure of set with monoid of operators \(\mathbf{O}(X')\), in a
manner functorial in \(X'\). For this one constructs the following diagram,
where \(X_{X'}\) denotes \(X\times_S X'\) and \(f'\) the section of
\(X_{X'}\) over \(X'\) defined by \(f\colon X'\to X\).
% typo?: kept as printed (diagram arrowheads: left-hand down, right-hand up)
\[
\xymatrix@C=1.6em@R=1.5em{
  & \mathrm{T}_{X_{X'}/X'}(\mathcal{M}) \ar[dl] & \\
  \mathrm{T}_{X/S}(\mathcal{M})
    & & \mathrm{T}_{X/S}(\mathcal{M})_{X'} \ar[ul]
}
\]
\[
\xymatrix@C=2.2em@R=1.5em{
  & X_{X'} \ar[dl] & \\
  X \ar[dr]
    & & X' \ar[ul]_{f'} \ar[ll]^{f} \ar[dl] \\
  & S &
}
\]
{\renewcommand{\thefootnote}{(20)}\nde{The original has been set out in
detail in what follows.}}
This diagram, together with 3.2.1, shows that
\(\mathrm{T}_{X/S}(\mathcal{M})_{X'}(X')\) is identified with
\[
\mathrm{L}^{f'}_{X_{X'}/X'}(\mathcal{M})(X')
=\{\text{\(X'\)-morphisms }
\psi\colon\mathrm{I}_{X'}(\mathcal{M})\to X_{X'}
\text{ such that }\psi\circ\varepsilon_{\mathcal{M}}=f'\},
\]
on which every \(a\in\mathbf{O}(X')\) operates via its action on
\(\mathrm{I}_{X'}(\mathcal{M})\), \ie with the notations of~2.1.1, one has:
\(a\psi=\psi\circ a^*\), \ie for every \(X''\to X'\) and
\(x\in\mathrm{I}_{X'}(\mathcal{M})(X'')\),
\((a\psi)(x)=\psi(x\cdot a)\).
One then checks easily that this construction is functorial in \(X'\).

The isomorphisms of Proposition~3.4 are then by construction isomorphisms
for the structures of \(\mathbf{O}_{X_{S'}}\)-objects, \resp
\(\mathbf{O}_{S'}\)-objects.

\begin{remark}{3.4.2}
\label{II.3.4.2}
{\renewcommand{\thefootnote}{(21)}\nde{Remarks 3.4.2 and 3.4.3 have been
added.}}
The operation of \(\mathbf{O}_X\) on \(\mathrm{T}_{X/S}(\mathcal{M})\) may be
seen, more simply, as follows. For every \(f\colon X'\to X\), one has
\[
\Hom_X(X',\mathrm{T}_{X/S}(\mathcal{M}))
=\{\phi\in\Hom_S(\mathrm{I}_{X'}(\mathcal{M}),X)
\mid\phi\circ\varepsilon_{\mathcal{M}}=f\},
\]
and one has seen in~2.1.3 that \(\mathrm{I}_{X'}(\mathcal{M})\), considered as
an \(S\)-functor, is endowed with an operation of the monoid
\(\mathbf{O}(X')\) which preserves the zero section
\(\varepsilon_{\mathcal{M}}\colon X'\to\mathrm{I}_{X'}(\mathcal{M})\).
Consequently, if one denotes by \(a^*\) the \(X'\)-endomorphism of
\(\mathrm{I}_{X'}(\mathcal{M})\) defined by \(a\in\mathbf{O}(X')\), one has:
\(a\phi=\phi\circ a^*\), \ie for every \(S'\to S\) and
\((m,x')\in\Hom_S(S',\mathrm{I}_S(\mathcal{M})\times_S X')\),
\[
(a\phi)(m,x')=\phi(m\cdot a(x'),x')
\]
(note that \(a^*\circ\varepsilon_{\mathcal{M}}=\varepsilon_{\mathcal{M}}\),
whence \((a\phi)\circ\varepsilon_{\mathcal{M}}=\phi\circ\varepsilon_{\mathcal{M}}=f\)).

Likewise, the operation of \(\mathbf{O}_S\) on
\(\mathrm{L}^{u}_{X/S}(\mathcal{M})\) may be described as follows. For every
\(t\colon S'\to S\), \(\mathrm{L}^{u}_{X/S}(\mathcal{M})(S')\) is the set of
\(S\)-morphisms \(\phi\colon\mathrm{I}_{S'}(\mathcal{M})\to X\) such that
\(\phi\circ\varepsilon_{\mathcal{M}}=u\circ t\); for such a \(\phi\) and
\(a\in\mathbf{O}(S')\), one has: \(a\phi=\phi\circ a^*\).
\end{remark}

\begin{remark}{3.4.3}
\label{II.3.4.3}
\textup{\textsuperscript{(21)}}
The observations of~3.1.2 hold equally for the operation of
\(\mathbf{O}_S\) on \(\mathrm{L}^{u}_{X/S}(\mathcal{M})\) and for that of
\(\mathbf{O}_X\) on \(\mathrm{T}_{X/S}(\mathcal{M})\).
\end{remark}

\begin{definition}{3.5}
\label{II.3.5}
Let \(S\) be a scheme and \(X\) an \(S\)-functor. One says that \(X\) satisfies
the condition~(E) relative to \(S\) if, for every \(S'\) over \(S\) and all free
\(\OO_{S'}\)-modules of finite type \(\mathcal{M}\) and \(\mathcal{N}\), the
diagram of sets
\[
\xymatrix@C=1.4em@R=1.4em{
  & X(\mathrm{I}_{S'}(\mathcal{M}\oplus\mathcal{N})) \ar[dl] \ar[dr] & \\
  X(\mathrm{I}_{S'}(\mathcal{M})) \ar[dr]
    & & X(\mathrm{I}_{S'}(\mathcal{N})) \ar[dl] \\
  & X(S') &
}
\]
obtained by applying \(X\) to the diagram \((*)\) defined in~2.1, is
cartesian.%
{\renewcommand{\thefootnote}{(22)}\nde{For examples of \(S\)-functors and
\(S\)-groups not satisfying~(E), see 6.2 and~6.3.}}
\end{definition}

\numpara{3.5.1}
\label{II.3.5.1}
\emph{It comes to the same to say that the functor
\(\mathcal{M}\mapsto\mathrm{T}_{X/S}(\mathcal{M})\) takes direct sums of free
\(\OO_S\)-modules of finite type to products of \(X\)-functors};%
{\renewcommand{\thefootnote}{(23)}\nde{What follows has been added.}}
in this case, the same holds for the functor
\(\mathcal{M}\mapsto\mathrm{L}^{u}_{X/S}(\mathcal{M})
=S\times_X\mathrm{T}_{X/S}(\mathcal{M})\),
for every \(u\in\Gamma(X/S)\).

Proposition~2.2 shows that \emph{every representable functor satisfies the
condition~(E)}.

\emph{Abbreviation}: instead of saying ``\(X\) satisfies the condition~(E)
with respect to \(S\)'', one will sometimes say ``\(X/S\) satisfies the
condition~(E)''.

If \(X/S\) satisfies the condition~(E), the functor
\(\mathcal{M}\mapsto\mathrm{T}_{X/S}(\mathcal{M})\) commutes with products,
hence takes groups to groups. In particular
\(\mathrm{T}_{X/S}(\mathcal{M})\) is a commutative \(X\)-group. For the same
reason, \(\mathrm{L}^{u}_{X/S}(\mathcal{M})\) is a commutative \(S\)-group.

\begin{proposition}{3.6}
\label{II.3.6}
If \(X/S\) satisfies~(E), the structure of abelian group on
\(\mathrm{T}_{X/S}(\mathcal{M})\)
(\resp \(\mathrm{L}^{u}_{X/S}(\mathcal{M})\)) and the operation of
\(\mathbf{O}_X\) (\resp \(\mathbf{O}_S\)) endow
\(\mathrm{T}_{X/S}(\mathcal{M})\)
(\resp \(\mathrm{L}^{u}_{X/S}(\mathcal{M})\)) with a structure of
\(\mathbf{O}_X\)-module (\resp \(\mathbf{O}_S\)-module).
\end{proposition}

The operation of \(\mathbf{O}_X\) (\resp \(\mathbf{O}_S\)) is functorial in
\(\mathcal{M}\); it therefore respects the structure of abelian group which
is deduced by functoriality from that of \(\mathcal{M}\).%
{\renewcommand{\thefootnote}{(24)}\nde{What follows has been set out in
detail.}}
Indeed, let us take up again the notations of~3.1.2. The structure of
(abelian) \(X\)-group on \(\mathrm{T}_{X/S}(\mathcal{M})\) is defined by the
composite:
\[
\mathrm{T}_{X/S}(\mathcal{M})\times_X\mathrm{T}_{X/S}(\mathcal{M})
\simeq
\mathrm{T}_{X/S}(\mathcal{M}\oplus\mathcal{M})
\xrightarrow{m}
\mathrm{T}_{X/S}(\mathcal{M}),
\]
and on the other hand the morphism
\[
\mathrm{T}_{X/S}(\mathcal{M})
\xrightarrow{\delta}
\mathrm{T}_{X/S}(\mathcal{M}\oplus\mathcal{M})
\simeq
\mathrm{T}_{X/S}(\mathcal{M})\times_X\mathrm{T}_{X/S}(\mathcal{M})
\]
is the diagonal morphism. Taking Remark~3.4.3 into account, one deduces from
the equalities \((\dagger)\) of~3.1.2 that
\[
\lambda(x+y)=\lambda x+\lambda y,\qquad (\lambda+\mu)x=\lambda x+\mu x,
\]
for every \(f\colon X'\to X\),
\(x,y\in\Hom_X(X',\mathrm{T}_{X/S}(\mathcal{M}))\) and
\(\lambda,\mu\in\mathbf{O}(X')\).

\begin{remark}{3.6.1}
\label{II.3.6.1}
If \(X\) is representable, in which case, on the one hand it satisfies~(E),
and on the other hand \(\mathrm{T}_{X/S}\) and \(\mathrm{L}^{u}_{X/S}\) are
representable by vector fibrations, the preceding laws are the same as those
which are deduced from the structures of vector fibration
(\cf I~4.6).%
{\renewcommand{\thefootnote}{(25)}\nde{That is, for every \(S\)-morphism
\(f\colon X'\to X\), the action of \(\mathbf{O}(X')\) on
\(\Hom_X(X',\mathrm{T}_{X/S}(\mathcal{M}))\) corresponds, via the
identification
\(\Hom_X(X',\mathrm{T}_{X/S}(\mathcal{M}))
\simeq
\Hom_{\OO_{X'}}\bigl(f^*(\Omega^1_{X/S}),
\mathcal{M}\otimes_{\OO_S}\OO_{X'}\bigr)\),
to the natural action of \(\mathbf{O}(X')\) on the term on the right; this
follows from the proof of SGA~1, III~5.1 (see also the addition~0.1.8 in
Expos\'e~III).}}
\end{remark}

\begin{proposition}{3.4. bis}
\label{II.3.4.bis}
If \(X/S\) satisfies~(E), then \(X_{S'}/S'\) satisfies~(E) and the
isomorphisms of~3.4 respect the structures of
\(\mathbf{O}_{X_{S'}}\)-modules, \resp of \(\mathbf{O}_{S'}\)-modules.
\end{proposition}

Without comment.

\begin{proposition}{3.7}
\label{II.3.7}
The functors \(\mathrm{T}_{X/S}(\mathcal{M})\) and
\(\mathrm{L}^{u}_{X/S}(\mathcal{M})\) are functorial in \(X\), \ie if
\(f\colon X\to X'\) is an \(S\)-morphism, one has commutative diagrams:
\[
\xymatrix{
  \mathrm{T}_{X/S}(\mathcal{M}) \ar[r]^{\mathrm{T}(f)} \ar[d]
    & \mathrm{T}_{X'/S}(\mathcal{M}) \ar[d] \\
  X \ar[r]
    & X',
}
\qquad
\xymatrix{
  \mathrm{L}^{u}_{X/S}(\mathcal{M}) \ar[rr]^{\mathrm{L}(f)} \ar[dr]
    & & \mathrm{L}^{f\circ u}_{X'/S}(\mathcal{M}) \ar[dl] \\
  & S. &
}
\]
\end{proposition}

{\renewcommand{\thefootnote}{(26)}\nde{The sentence which follows has been
added.}}
\emph{Moreover, if \(f\) is a monomorphism, the same holds of
\(\mathrm{T}(f)\) and \(\mathrm{L}(f)\).}

The existence of \(\mathrm{T}(f)\) and \(\mathrm{L}(f)\), as well as the last
assertion, are deduced immediately from the definitions. The commutativity of
the diagrams then results from the functoriality of these morphisms with
respect to \(\mathcal{M}\) and from the fact that
\(\mathrm{T}_{X/S}(0)=X\).

\begin{remark}{3.7.1}
\label{II.3.7.1}
{\renewcommand{\thefootnote}{(27)}\nde{The hypothesis that \(X\) and \(X'\)
are representable has been added, and what follows has been set out in
detail.}}
Suppose \(X\) and \(X'\) representable, and let \(r\) be the rank of the free
\(\OO_S\)-module \(\mathcal{M}\). Then, by~2.2.2,
\(\mathrm{T}_{X/S}(\mathcal{M})\) is isomorphic to the product over \(X\) of
\(r\) copies of \(\mathbf{V}(\Omega^1_{X/S})\), and likewise for
\(\mathrm{T}_{X'/S}(\mathcal{M})\). Consequently, the square above is
cartesian when \(f\) is an open immersion, more generally when
\(f^*(\Omega^1_{X'/S})=\Omega^1_{X/S}\), for example if \(f\) is \'etale;
under these conditions, one has an isomorphism of \(S\)-functors
\[
\mathrm{L}^{u}_{X/S}(\mathcal{M})
\isomto
\mathrm{L}^{f\circ u}_{X'/S}(\mathcal{M}).
\]
\end{remark}
